\documentclass[11pt]{article}

\usepackage[utf8]{inputenc}
\usepackage[T1]{fontenc}
\usepackage{amsmath}
\usepackage{amssymb}
\usepackage{booktabs}
\usepackage{geometry}
\usepackage{longtable}
\usepackage{array}
\usepackage{xcolor}
\usepackage{hyperref}
\usepackage{enumitem}

\geometry{margin=1in}
\hypersetup{
  colorlinks=true,
  linkcolor=blue!55!black,
  urlcolor=blue!55!black,
  citecolor=blue!55!black,
  pdftitle={Next Analysis Plan: Variance, PCA, Low-Rank Variants, and Feature Decisions},
  pdfauthor={TFM Bayesian Football Analytics Pipeline}
}

\setlength{\parindent}{0pt}
\setlength{\parskip}{0.55em}
\renewcommand{\arraystretch}{1.15}
\setlength{\tabcolsep}{4pt}
\emergencystretch=2em
\newcolumntype{L}[1]{>{\raggedright\arraybackslash}p{#1}}
\newcolumntype{R}[1]{>{\raggedleft\arraybackslash}p{#1}}
\newcommand{\code}[1]{\texttt{\detokenize{#1}}}

\title{\textbf{Next Analysis Plan}\\
\large Variance Decomposition, Diagonal Projections, PCA, Low-Rank Variants, Heatmaps, and Feature Decisions}
\author{Bayesian Football Analytics Pipeline}
\date{Planned from current fitted results}

\begin{document}
\maketitle

\begin{abstract}
This document is a planning note only. It does not propose immediate code
changes. The goal is to define the next analytical work needed after the
current real-data diagonal fit, simulation recovery checks, and validation
outputs. The plan focuses on variance decomposition, projections using
\(\Sigma_a\) and \(\Sigma_b\), PCA-style summaries, additional low-rank
simulation variants, sign-aware heatmaps of simulated effects, residual-sorted
explained-variance summaries, and a decision framework for retaining or
excluding \code{rate_*} variables.
\end{abstract}

\section{Current Starting Point}

The current validated pipeline has three main fitted objects:

\begin{itemize}[leftmargin=*]
  \item real-data diagonal additive model;
  \item simulated-data diagonal additive model fitted to low-rank simulated data;
  \item simulated-data low-rank recovery model.
\end{itemize}

The main existing conclusion is that the diagonal additive model is suitable
for descriptive variance decomposition, while the current full low-rank
recovery model has poor loading-level convergence. The next work should
therefore separate two goals:

\begin{enumerate}[leftmargin=*]
  \item strengthen the thesis-facing interpretation of the diagonal model;
  \item test narrower low-rank variants to understand whether the failure is
  caused by player covariance, season covariance, their combination, or limited
  season information.
\end{enumerate}

\section{Main Questions To Answer}

\begin{longtable}{L{0.27\textwidth}L{0.37\textwidth}L{0.27\textwidth}}
\toprule
Question & Why it matters & Best summary \\
\midrule
\endhead
Which features are mostly player, season, or residual driven? &
This is the main thesis interpretation of the additive model. &
Variance-share table and stacked bar plot sorted by residual share. \\
\addlinespace
What do diagonal \(\Sigma_a\) and \(\Sigma_b\) imply geometrically? &
Diagonal covariance matrices do not create rotated latent dimensions, but they
still define feature-wise axes of player and season variation. &
Feature ranking, diagonal projection scores, and player/season dominance maps. \\
\addlinespace
Does PCA add information beyond diagonal variance ranking? &
For diagonal matrices, PCA is trivial; for empirical effects or low-rank
covariances, PCA can reveal correlated feature profiles. &
Compare PCA on empirical/posterior effects and on low-rank covariance estimates. \\
\addlinespace
Where exactly does low-rank recovery fail? &
The current low-rank model mixes poorly; isolated experiments can diagnose
whether the issue is player block, season block, or both. &
Simulation recovery table across model variants. \\
\addlinespace
How should signs of simulated player and season effects be shown? &
Heatmaps can be misleading unless factor signs are aligned and rows are ordered
sensibly. &
Sign-aligned heatmaps for loadings/effects plus clustered summaries. \\
\addlinespace
Should \code{rate_*} variables remain in the model? &
Many rate variables are residual dominated, potentially noisy, and denominator
sensitive. &
Rate-variable audit and optional sensitivity comparison with rates removed. \\
\bottomrule
\end{longtable}

\section{Work Package 1: Variance Decomposition}

\subsection{Objective}

Create a thesis-ready variance-decomposition summary for every fitted model
where variance components are available. For each feature \(p\), compute:

\begin{align*}
\text{player share}_p &=
\frac{\sigma^2_{a,p}}{\sigma^2_{a,p} + \sigma^2_{b,p} + \sigma^2_{e,p}}, \\
\text{season share}_p &=
\frac{\sigma^2_{b,p}}{\sigma^2_{a,p} + \sigma^2_{b,p} + \sigma^2_{e,p}}, \\
\text{residual share}_p &=
\frac{\sigma^2_{e,p}}{\sigma^2_{a,p} + \sigma^2_{b,p} + \sigma^2_{e,p}}.
\end{align*}

\subsection{Recommended Summaries}

\begin{enumerate}[leftmargin=*]
  \item \textbf{All-feature variance table}: one row per feature with posterior
  mean player, season, residual, total variance, and the three shares.
  \item \textbf{Residual-sorted table}: same table sorted descending by residual
  share. This directly identifies features least explained by stable player or
  season effects.
  \item \textbf{Player-sorted and season-sorted tables}: top 15 features by
  player share and top 15 by season share.
  \item \textbf{Stacked bar plot}: one horizontal bar per feature, sorted by
  residual share, with player, season, and residual colors.
  \item \textbf{Feature-family summary}: compare \code{per90_*} versus
  \code{rate_*} variables using median and interquartile range of player,
  season, and residual shares.
\end{enumerate}

\subsection{Interpretation Rules}

\begin{longtable}{L{0.30\textwidth}L{0.60\textwidth}}
\toprule
Pattern & Interpretation \\
\midrule
\endhead
High player share & The feature is stable across observed seasons for players;
it is suitable for player-profile interpretation. \\
High season share & The feature is sensitive to common season-level shifts,
possibly tactical, league-wide, measurement, or era effects. \\
High residual share & The feature is noisy or situational at player-season
level after accounting for player and season. \\
High residual share for rates & Check denominator instability, rare events,
and whether the rate should be retained. \\
\bottomrule
\end{longtable}

\section{Work Package 2: Projections Using Diagonal \texorpdfstring{\(\Sigma_a\) And \(\Sigma_b\)}{Sigma a And Sigma b}}

\subsection{Key Point}

For the diagonal model, \(\Sigma_a\) and \(\Sigma_b\) are diagonal matrices.
Therefore they do not define rotated multivariate directions. Their eigenvectors
are the original features, and their eigenvalues are the feature-specific
variances. The right interpretation is not a latent PCA rotation, but a
feature-axis weighting.

\subsection{Planned Projections}

\begin{enumerate}[leftmargin=*]
  \item \textbf{Player-variance projection}: weight each feature by
  \(\sigma^2_{a,p}\) or by player share. This highlights the part of each
  observation most aligned with stable player variation.
  \item \textbf{Season-variance projection}: weight each feature by
  \(\sigma^2_{b,p}\) or by season share. This highlights features most aligned
  with season-level shifts.
  \item \textbf{Residual projection}: weight each feature by residual share.
  This highlights noisy or weakly explained features.
  \item \textbf{Player-versus-season contrast}: compute a feature contrast such
  as \(\log(\sigma^2_{a,p} / \sigma^2_{b,p})\), with a small stabilization
  constant if needed. This ranks features by whether they are more player-like
  or season-like.
\end{enumerate}

\subsection{Recommended Outputs}

\begin{longtable}{L{0.32\textwidth}L{0.58\textwidth}}
\toprule
Output & Purpose \\
\midrule
\endhead
Diagonal variance scree plots for \(\Sigma_a\) and \(\Sigma_b\) &
Show whether a few features dominate player or season variance. \\
Feature-axis projection table &
Rank features by player, season, residual, and player-season contrast. \\
Two-dimensional map: player share vs season share &
Identify player-dominant, season-dominant, residual-dominant, and mixed
features. \\
Feature-family overlay &
Show whether \code{rate_*} variables cluster in the residual-dominant area. \\
\bottomrule
\end{longtable}

\section{Work Package 3: PCA}

\subsection{What PCA Should And Should Not Mean Here}

PCA on a diagonal covariance matrix is mostly a sorted variance ranking. It
does not reveal cross-feature structure because off-diagonal covariance is zero.
Therefore PCA should be used in three more informative places:

\begin{enumerate}[leftmargin=*]
  \item \textbf{PCA on posterior player effects \(a_i\)} from the diagonal
  model. This asks whether estimated player profiles vary along interpretable
  multifeature dimensions.
  \item \textbf{PCA on posterior season effects \(b_j\)} from the diagonal
  model. This asks whether seasons separate along common tactical or
  measurement patterns.
  \item \textbf{PCA/eigendecomposition of low-rank covariance estimates} in
  models where \(\Sigma_a = \Lambda_a\Lambda_a^\top + \Psi_a\) or
  \(\Sigma_b = \Lambda_b\Lambda_b^\top + \Psi_b\). This asks whether the
  low-rank structure creates interpretable shared dimensions.
\end{enumerate}

\subsection{Recommended PCA Summaries}

\begin{longtable}{L{0.30\textwidth}L{0.60\textwidth}}
\toprule
Summary & Why it makes sense \\
\midrule
\endhead
Explained variance ratio by PC &
Shows whether player or season effects have low-dimensional structure. \\
Top positive and negative loadings per PC &
Makes each PC interpretable in football terms. \\
PC score scatterplots for players &
Shows whether player profiles occupy meaningful regions; optionally label only
extreme players to avoid clutter. \\
PC score plot for seasons &
Only 12 seasons exist, so all seasons can be labelled. \\
Bootstrap/posterior uncertainty for loadings &
Prevents overinterpreting unstable PCA directions. \\
\bottomrule
\end{longtable}

\section{Work Package 4: New Low-Rank Simulation Variants}

\subsection{Objective}

The current low-rank recovery model is too broad to diagnose the failure
precisely. The next experiments should isolate player and season covariance
complexity.

\subsection{Recommended Experiment Order}

\begin{longtable}{L{0.20\textwidth}L{0.26\textwidth}L{0.43\textwidth}}
\toprule
Experiment & Model structure & Reason \\
\midrule
\endhead
LR-A0 & \(a_i\) low-rank, no season effect \(b_j\) &
Start with the simplest player-only low-rank recovery. This checks whether
player low-rank structure is identifiable without season competition. \\
\addlinespace
LR-A1 & \(a_i\) low-rank, \(b_j\) diagonal &
Adds season adjustment while keeping season covariance simple. This tests
whether player low-rank recovery survives realistic season controls. \\
\addlinespace
LR-B1 & \(a_i\) diagonal, \(b_j\) low-rank &
Tests whether season low-rank structure is recoverable at all with only 12
seasons. This is expected to be difficult. \\
\addlinespace
LR-AB & \(a_i\) low-rank, \(b_j\) low-rank &
Current richest structure. Revisit only after the simpler cases are understood. \\
\bottomrule
\end{longtable}

\subsection{Recovery Metrics For Every Variant}

Each simulation variant should produce the same recovery table:

\begin{itemize}[leftmargin=*]
  \item sampler health: divergences, treedepth hits, E-BFMI, acceptance;
  \item posterior health: max Rhat, number of Rhat warnings, min ESS;
  \item loading recovery: bias, MAE, RMSE, correlation, 90 percent coverage;
  \item covariance diagonal recovery: same metrics for \(\Sigma_a\) and/or
  \(\Sigma_b\);
  \item residual scale recovery;
  \item feature-level recovery sorted by worst absolute error.
\end{itemize}

\subsection{Decision Logic}

\begin{longtable}{L{0.36\textwidth}L{0.54\textwidth}}
\toprule
Outcome & Interpretation \\
\midrule
\endhead
LR-A0 works but LR-A1 fails & Season adjustment interferes with player
low-rank identification. \\
LR-A1 works but LR-B1 fails & Player low-rank is feasible; season low-rank is
not reliable with 12 seasons. \\
LR-A0 fails & The player low-rank specification or constraints need stronger
identifiability before adding seasons. \\
LR-B1 fails but diagonal season works & Keep season diagonal in thesis models. \\
LR-AB works only with strong anchors or truth initialization & Treat low-rank
as a sensitivity/recovery experiment, not as the main empirical model. \\
\bottomrule
\end{longtable}

\section{Work Package 5: Sign-Aware Heatmaps For Simulated \texorpdfstring{\(a_i\) And \(b_j\)}{a i And b j}}

\subsection{Issue}

Low-rank factor signs are not intrinsically identified. A factor can change
sign if its loading and score both change sign. Therefore heatmaps of simulated
or recovered factors need sign alignment before interpretation.

\subsection{Recommended Heatmaps}

\begin{enumerate}[leftmargin=*]
  \item \textbf{Truth heatmaps}: simulated true \(a_i\) and \(b_j\) effects by
  feature. These show the actual generated signal.
  \item \textbf{Posterior mean heatmaps}: recovered posterior mean effects,
  aligned to truth where simulation truth exists.
  \item \textbf{Error heatmaps}: posterior mean minus truth, using a diverging
  color scale centered at zero.
  \item \textbf{Sign agreement heatmaps}: sign of posterior mean compared with
  sign of truth, useful when magnitude recovery and sign recovery differ.
  \item \textbf{Uncertainty heatmaps}: posterior probability that an effect is
  positive, \(P(a_{ip} > 0)\) or \(P(b_{jp} > 0)\), where available.
\end{enumerate}

\subsection{Display Strategy}

Showing every player and every feature in one figure may be too dense because
there are 1,650 players and 53 features. The plan should include:

\begin{itemize}[leftmargin=*]
  \item full heatmap as an appendix artifact;
  \item clustered heatmap by player-effect profile;
  \item top 50 players by absolute total effect norm;
  \item repeated-player-only heatmap, because repeated observations identify
  player effects better;
  \item all 12 seasons labelled individually, because the season heatmap is
  small enough to read.
\end{itemize}

\section{Work Package 6: Explained Variance Sorted By Residual}

\subsection{Objective}

For every completed model where variance components are meaningful, produce a
single comparable explained-variance table sorted by residual share.

\subsection{Recommended Master Table}

\begin{longtable}{L{0.27\textwidth}L{0.10\textwidth}R{0.08\textwidth}R{0.08\textwidth}R{0.10\textwidth}R{0.07\textwidth}L{0.09\textwidth}}
\toprule
Feature & Family & Player & Season & Residual & Total & Flag \\
\midrule
\endhead
\small rate\_successful\_crosses & rate & -- & -- & high & -- & review \\
\small rate\_successful\_sliding\newline \small \_tackles & rate & -- & -- & high & -- & review \\
\small rate\_shots\_on\_target & rate & -- & -- & high & -- & review \\
\small rate\_goal\_conversion & rate & -- & -- & high & -- & review \\
\bottomrule
\end{longtable}

The actual table should contain all 53 features. The rows above illustrate the
expected high-residual rate-variable cases from the current fitted results.

\subsection{Useful Flags}

\begin{itemize}[leftmargin=*]
  \item \textbf{player-dominant}: player share above a chosen threshold;
  \item \textbf{season-dominant}: season share above a chosen threshold;
  \item \textbf{residual-dominant}: residual share above a chosen threshold;
  \item \textbf{rate-review}: variable name starts with \code{rate_*} and has
  high residual share;
  \item \textbf{low-ESS}: variance component has weak posterior effective
  sample size;
  \item \textbf{Rhat-warning}: variance component has convergence warning.
\end{itemize}

\section{Work Package 7: Decide Whether \code{rate_*} Variables Are Needed}

\subsection{Reason To Reconsider Rates}

Rate variables can be statistically fragile because they depend on
denominators. A player with few attempts can have an extreme rate that is not a
stable trait. The current real-data model already suggests that several rate
features are mostly residual variation.

\subsection{Decision Framework}

\begin{longtable}{L{0.30\textwidth}L{0.60\textwidth}}
\toprule
Criterion & Action \\
\midrule
\endhead
High residual share and rare denominator & Strong candidate for removal or
separate modelling. \\
High player share and adequate denominator & Candidate to keep as a stable
skill-like feature. \\
High correlation with a cleaner \code{per90_*} variable & Consider removing
the noisier rate variable. \\
Important football interpretation despite noise & Keep, but report caveat. \\
Unstable under sensitivity analysis & Remove from main model or move to
appendix. \\
\bottomrule
\end{longtable}

\subsection{Recommended Sensitivity Comparison}

Do not remove all rates automatically. Instead, plan one sensitivity fit or
analysis with:

\begin{itemize}[leftmargin=*]
  \item full feature set, current baseline;
  \item no \code{rate_*} variables;
  \item optional filtered-rate set, keeping only rates with adequate
  denominators and non-extreme residual share.
\end{itemize}

Compare these on:

\begin{itemize}[leftmargin=*]
  \item stability of top player-dominant \code{per90_*} features;
  \item posterior health;
  \item residual-share distribution;
  \item whether season-level conclusions change.
\end{itemize}

\section{Prioritized Execution Plan}

\begin{longtable}{L{0.12\textwidth}L{0.54\textwidth}L{0.24\textwidth}}
\toprule
Priority & Task & Output \\
\midrule
\endhead
1 & Create final variance-decomposition tables sorted by residual, player, and
season share. & Thesis-ready CSV/table and stacked bar plot. \\
2 & Add feature-family summaries for \code{per90_*} vs \code{rate_*}. &
Rate-variable audit table. \\
3 & Produce diagonal \(\Sigma_a\), \(\Sigma_b\), and residual projection
summaries. & Feature-axis ranking and player-season contrast plot. \\
4 & Run PCA summaries on posterior player and season effects, not merely on
diagonal covariance matrices. & PC explained variance, loading tables, score
plots. \\
5 & Design and run low-rank simulation variants in order: LR-A0, LR-A1, LR-B1,
then LR-AB only if justified. & Recovery comparison table. \\
6 & Produce sign-aligned heatmaps for simulated truth, posterior means, errors,
and sign agreement. & Heatmap appendix and selected readable panels. \\
7 & Decide whether \code{rate_*} variables stay in the main thesis model. &
Feature inclusion recommendation. \\
\bottomrule
\end{longtable}

\section{Recommended Final Summary Figures}

\begin{enumerate}[leftmargin=*]
  \item \textbf{Variance decomposition by feature}, sorted by residual share.
  This should be the main thesis figure.
  \item \textbf{Player share vs season share scatter}, with point shape or color
  for \code{per90_*} versus \code{rate_*}.
  \item \textbf{Residual-share distribution by feature family}, comparing
  \code{per90_*} and \code{rate_*}.
  \item \textbf{PCA loading table and score plot for player effects}, if PCs are
  stable.
  \item \textbf{Season-effect heatmap}, all 12 seasons by selected features.
  \item \textbf{Low-rank recovery dashboard}, one table comparing LR-A0, LR-A1,
  LR-B1, and LR-AB on convergence and recovery metrics.
  \item \textbf{Simulation heatmap appendix}, truth, posterior mean, and error
  maps after sign alignment.
\end{enumerate}

\section{Expected Interpretation If Current Patterns Persist}

If the current results remain stable, the likely final story is:

\begin{itemize}[leftmargin=*]
  \item the diagonal additive model is the main empirical model because it gives
  interpretable and reasonably stable variance decomposition;
  \item many \code{per90_*} volume and intensity variables are stable
  player-profile measures;
  \item several \code{rate_*} variables are mostly residual and should be
  treated cautiously;
  \item season-level effects exist for selected tactical or measurement-sensitive
  features, but broad season covariance is hard to estimate with only 12 seasons;
  \item low-rank player covariance may be feasible in simplified simulations,
  but low-rank season covariance is expected to be fragile;
  \item full low-rank player-plus-season recovery should not be the main thesis
  claim unless the simpler variants show reliable convergence and recovery.
\end{itemize}

\section{Immediate Non-Coding Next Step}

Before implementing any new model variant, define the final reporting tables
and figures exactly. The first deliverable should be a compact specification
listing each output, its columns, sorting rule, and interpretation sentence.
That avoids running additional models before the thesis summaries are clear.

\end{document}
